% Compilation (depuis la racine du dépôt) : bash tools/build_report.sh
\documentclass[11pt,a4paper]{article}

\usepackage{fontspec}
\setmonofont{DejaVu Sans Mono}[Scale=0.88]
\usepackage[french]{babel}
\usepackage[a4paper,margin=2.2cm]{geometry}
\usepackage{graphicx,float,xcolor,fancyhdr,caption}
\usepackage[hidelinks]{hyperref}
\graphicspath{{./}}

\setlength{\parindent}{0pt}
\setlength{\parskip}{0.6em}
\captionsetup{font=small,labelfont=bf}

\pagestyle{fancy}
\fancyhf{}
\lhead{\small TP Sécurité Oracle -- RBAC}
\rhead{\small LACHHAB Mohammed -- GINF2}
\cfoot{\thepage}

% \code{R_RH_LECTEUR} : identifiants, codes ORA-, commandes (les _ # % sont gérés)
\newcommand{\code}[1]{\texttt{\detokenize{#1}}}
% \shots{1_1} : inclut les images produites par tools/shot.py (screenshots/1_1.tex)
\newcommand{\shots}[1]{\IfFileExists{screenshots/#1.tex}{\input{screenshots/#1.tex}}{\textbf{[capture manquante : #1]}}}
% \ecart{texte} : encadré « écart 19c / 26ai »
\newcommand{\ecart}[1]{\par\noindent\fbox{\parbox{\dimexpr\linewidth-2\fboxsep-2\fboxrule\relax}{\small\textbf{Écart par rapport au TP (19c $\to$ 26ai) :} #1}}\par}

\begin{document}

% ------------------------------------------------------------------ PAGE DE GARDE
\begin{titlepage}
\centering
\includegraphics[width=0.32\textwidth]{Ensat.jpg}\par
\vspace{0.8cm}
{\large\bfseries École Nationale des Sciences Appliquées de Tanger\par}
\vspace{0.2cm}
{\normalsize Université Abdelmalek Essaâdi\par}
\vspace{2.2cm}
\rule{\textwidth}{0.8pt}\par
\vspace{0.5cm}
{\Huge\bfseries Compte rendu de travaux pratiques\par}
\vspace{0.6cm}
{\LARGE Sécurité : contrôle d'accès basé sur les rôles (RBAC)\\[0.2cm] avec Oracle Database\par}
\vspace{0.5cm}
\rule{\textwidth}{0.8pt}\par
\vspace{1.4cm}
{\large Module : Administration Oracle\par}
\vfill
\begin{tabular}{rl}
  \textbf{Réalisé par :}        & LACHHAB Mohammed \\[0.2cm]
  \textbf{Filière :}            & GINF2 \\[0.2cm]
  \textbf{Encadrant :}          & Mr. BADIR Hassan \\[0.2cm]
  \textbf{Année universitaire :} & 2026/2027
\end{tabular}
\vfill
{\today}
\end{titlepage}

\tableofcontents
\newpage

% ------------------------------------------------------------------ ENVIRONNEMENT
\section*{Environnement de travail et conventions}
\addcontentsline{toc}{section}{Environnement de travail et conventions}

Le TP est rédigé pour Oracle Database 19c (CDB \code{ORCL}, PDB \code{ORCLPDB}). Il a été réalisé
sur \textbf{Oracle AI Database 26ai Free (version 23.26.3.0.0)}, exécutée dans un conteneur Docker
(image \code{gvenzl/oracle-free:slim}) hébergé dans un GitHub Codespace. Le conteneur est
une base multitenant (CDB \code{FREE}) ; la PDB utilisée est \textbf{\code{FREEPDB1}}, qui remplace
\code{ORCLPDB} partout dans l'énoncé. Toutes les connexions se font dans la PDB, sauf l'exercice 1.5
(racine \code{CDB\$ROOT}).

Les sessions sont exécutées avec SQL*Plus. Les « captures » de ce rapport sont des rendus d'images
réalisés à partir des journaux (\emph{logs}) réels des sessions ; le nom du journal est rappelé en pied
de chaque image. Quand le résultat obtenu diffère de ce que l'énoncé (version 19c) annonce, un encadré
\emph{Écart par rapport au TP} l'indique et l'explique.

\newpage

% ------------------------------------------------------------------ CHAPITRES
% Chaque fichier report/chN.tex commence par \section{...}. Les fichiers absents sont ignorés.
\IfFileExists{report/ch0.tex}{\section{Préparation}

\subsection{Test de l'outillage}
\textbf{Objectif.} Vérifier que SQL*Plus est opérationnel, que la session est bien ouverte dans la PDB
\code{FREEPDB1} et que le tablespace \code{USERS} existe avant de commencer le TP.

\shots{0_test}

\textbf{Observations.}
\begin{enumerate}
  \item Le script s'exécute via \code{bash tools/sq.sh 0_test < scripts/0\_test.sql}. Le journal
        \code{logs/0\_test.log} contient bien les commandes echoées (préfixe \code{SQL>}) suivies de
        leurs résultats : \code{-s} a été retiré de \code{tools/sq.sh} pour obtenir cette trace.
  \item \code{SHOW CON\_NAME} affiche \code{FREEPDB1} et \code{SELECT USER FROM dual} renvoie \code{SYS} :
        la connexion \code{AS SYSDBA} se fait bien dans la PDB et non dans \code{CDB\$ROOT}.
  \item \code{SELECT tablespace\_name FROM dba\_tablespaces} renvoie 5 tablespaces :
        \code{SYSTEM}, \code{SYSAUX}, \code{UNDOTBS1}, \code{TEMP} et \code{USERS}. Le tablespace
        \code{USERS} nécessaire au schéma RH est donc présent.
\end{enumerate}

\subsection{Exercice 0.1 -- Vérifier l'environnement}
\textbf{Objectif.} Charger le schéma RH puis contrôler son état (volumétrie, type d'authentification,
origine des comptes de la PDB).

Le script \code{scripts/00-init-rh.sql} (id \code{0\_init}) a été exécuté au préalable : il crée
l'utilisateur schéma-seul \code{RH} (\code{NO AUTHENTICATION}), les tables \code{departements} et
\code{employes}, insère 4 départements et 8 employés, puis crée la vue \code{v\_annuaire}. Toutes les
instructions ont abouti (\code{User created.}, \code{Table created.}, \code{View created.}).

\shots{0_1}

\textbf{Observations et réponses.}
\begin{enumerate}
  \item Nombre de lignes : \code{departements} = \textbf{4} et \code{employes} = \textbf{8}, conformément
        aux insertions du script d'initialisation.
  \item Pour le compte \code{RH}, \code{DBA\_USERS} renvoie : \code{AUTHENTICATION\_TYPE = NONE},
        \code{ACCOUNT\_STATUS = OPEN} et \code{ORACLE\_MAINTAINED = N}. Le type d'authentification vaut
        donc \textbf{\code{NONE}} : c'est un compte « schéma seul » (schema-only account), sans mot de
        passe, qui ne peut pas ouvrir de session ; il ne sert qu'à posséder des objets et n'est
        accessible qu'à travers un proxy ou un utilisateur applicatif.
  \item La PDB contient \textbf{28} utilisateurs maintenus par Oracle (\code{ORACLE\_MAINTAINED = 'Y'})
        et \textbf{2} utilisateurs non maintenus (\code{'N'}) : \code{RH}, créé par nous pour ce TP, et
        \code{PDBADMIN}, créé par le script de création de la PDB de l'image Docker. Ce dernier n'ayant
        pas été créé par Oracle lui-même, il est compté à part.
\end{enumerate}

\ecart{Le nombre d'utilisateurs maintenus par Oracle dépend de la version et de l'image. En 26ai Free
(23.26.3.0.0) on obtient 28 comptes \code{'Y'} ; ce total diffère de celui de l'énoncé 19c. Le compte
\code{PDBADMIN} livré par l'image explique le second utilisateur non maintenu.}
}{}
\IfFileExists{report/ch1.tex}{\section{Utilisateurs}

\subsection{Exercice 1.1 -- Premier compte}
\textbf{Objectif.} Créer un utilisateur sans aucun droit, constater qu'il ne peut pas se connecter, puis
lui accorder uniquement le privilège minimal.

\shots{1_1}

\textbf{Observations et réponses.}
\begin{enumerate}
  \item La création de \code{ALI} (\code{DEFAULT TABLESPACE users}, \code{TEMPORARY TABLESPACE temp})
        réussit : un utilisateur Oracle n'est qu'une identité, sans privilège implicite.
  \item La connexion \code{ali/"Ali\#2026"} échoue avec \code{ORA-01045: Login denied. User ALI does not
        have CREATE SESSION privilege}. Sans \code{CREATE SESSION}, aucun utilisateur ne peut ouvrir de
        session, même avec le bon mot de passe.
  \item Le privilège minimal est \code{GRANT CREATE SESSION TO ali}; après reconnexion,
        \code{SELECT USER FROM dual} renvoie \code{ALI}.
  \item \code{DBA\_USERS} affiche : \code{USERNAME = ALI}, \code{DEFAULT\_TABLESPACE = USERS},
        \code{TEMPORARY\_TABLESPACE = TEMP}, \code{PROFILE = DEFAULT}, \code{ACCOUNT\_STATUS = OPEN} et
        \code{CREATED = 29-OCT-26}.
\end{enumerate}

\subsection{Exercice 1.2 -- Quotas et création différée}
\textbf{Objectif.} Comprendre pourquoi une table peut être créée sans quota alors que l'insertion échoue,
puis observer le dépassement de quota.

\shots{1_2}

\textbf{Observations et réponses.}
\begin{enumerate}
  \item Après \code{GRANT CREATE TABLE TO ali}, ALI crée \code{notes(id NUMBER, texte VARCHAR2(4000))} :
        la commande réussit. L'insertion échoue avec \code{ORA-01950: The user ALI has insufficient quota
        on tablespace USERS}. La création avait réussi car l'allocation du segment est \emph{différée} au
        premier insert (segment deferred) : le dictionnaire enregistre la définition de la table, mais aucun
        espace n'est encore consommé.
  \item Après \code{ALTER USER ali QUOTA 1M ON users}, l'insertion réussit : \code{nb\_lignes = 1}.
  \item Le bloc PL/SQL de 2000 lignes s'arrête sur \code{ORA-01536: space quota exceeded for tablespace
        'USERS'} (remonté par \code{ORA-06512: at line 3}) : le mégaoctet concédé est saturé.
  \item \code{DBA\_TS\_QUOTAS} renvoie pour ALI sur \code{USERS} : \code{BYTES = 1048576} et
        \code{MAX\_BYTES = 1048576}. Un \code{MAX\_BYTES = -1} signifierait un quota \textbf{illimité}
        (la valeur \code{-1} code l'infini dans cette vue).
\end{enumerate}

\subsection{Exercice 1.3 -- Cycle de vie d'un compte}
\textbf{Objectif.} Expirer un mot de passe, le changer à la connexion, verrouiller puis déverrouiller le
compte et suivre les dates associées.

\shots{1_3}
\shots{1_3_chg}
\shots{1_3_lock}

\textbf{Observations et réponses.}
\begin{enumerate}
  \item \code{ALTER USER ali PASSWORD EXPIRE} passe \code{ACCOUNT\_STATUS} de \code{OPEN} à
        \code{EXPIRED} et renseigne \code{EXPIRY\_DATE = 29-OCT-26}. À la connexion, SQL*Plus affiche
        \code{ORA-28001: The account has expired. The password must be changed} puis propose
        \code{Changing password for ali / New password: / Retype new password:}. Le mot de passe est
        effectivement changé en \code{Ali\#2027} (message \code{Password changed}), puis le compte repasse
        \code{OPEN}.
  \item \code{ALTER USER ali ACCOUNT LOCK} passe l'état à \code{LOCKED} avec
        \code{LOCK\_DATE = 29-OCT-26}. La connexion renvoie alors \code{ORA-28000: The account is locked;
        login denied}.
  \item \code{ALTER USER ali ACCOUNT UNLOCK} ramène le compte à \code{OPEN} et vide \code{LOCK\_DATE}.
  \item Il est préférable de \emph{verrouiller} plutôt que de supprimer immédiatement le compte d'un
        employé qui part : on conserve ses objets (tables, vues), l'historique et la trace d'audit de ses
        actions, et l'opération est réversible. Une suppression \code{DROP USER ... CASCADE} détruit au
        contraire les objets et rend toute enquête impossible.
\end{enumerate}

\subsection{Exercice 1.4 -- Compte schéma seul et proxy}
\textbf{Objectif.} Manipuler un compte \code{NO AUTHENTICATION} et une connexion proxy.

\shots{1_4}

\textbf{Observations et réponses.}
\begin{enumerate}
  \item La connexion directe \code{rh/"nimportequoi"} échoue avec \code{ORA-01017: invalid credential or
        not authorized; logon denied} : un compte schéma seul n'a pas de mot de passe, donc aucune
        authentification directe n'est possible.
  \item Après \code{CREATE USER dev1 IDENTIFIED BY "Dev1\#2026"}, \code{GRANT CREATE SESSION TO dev1} et
        \code{ALTER USER rh GRANT CONNECT THROUGH dev1}, la première connexion
        \code{dev1[rh]/"Dev1\#2026"} échoue avec \code{ORA-01045: Login denied. User RH does not have
        CREATE SESSION privilege}. C'est bien \textbf{\code{RH}} qui manque du privilège : le proxy ouvre
        une session \emph{en tant que} RH, qui doit donc lui aussi posséder \code{CREATE SESSION}.
  \item Après \code{GRANT CREATE SESSION TO rh}, la connexion proxy réussit : \code{USER = RH} et
        \code{SYS\_CONTEXT('USERENV','PROXY\_USER') = DEV1}. Les tables du schéma courant
        (\code{user\_tables}) sont \code{DEPARTEMENTS} et \code{EMPLOYES}.
  \item Intérêt : le schéma \code{RH} n'a pas de mot de passe à voler ; chaque développeur s'authentifie
        avec son propre compte ; on peut retirer l'accès d'un seul développeur sans toucher au schéma ; et
        l'audit enregistre \code{PROXY\_USER} (DEV1) en plus de \code{USER} (RH), ce qui assure la
        traçabilité.
\end{enumerate}

\subsection{Exercice 1.5 -- Utilisateurs communs et locaux}
\textbf{Objectif.} Distinguer un utilisateur local d'un utilisateur commun (\code{CONTAINER=ALL}).

\shots{1_5}

\textbf{Observations et réponses.}
\begin{enumerate}
  \item Connecté à la racine (\code{SHOW CON\_NAME = CDB\$ROOT}), \code{CREATE USER bob IDENTIFIED BY
        "Bob\#2026"} \textbf{réussit} dans cet environnement (voir l'écart ci-dessous).
  \item \code{CREATE USER c##superviseur ... CONTAINER=ALL} puis \code{GRANT CREATE SESSION ...
        CONTAINER=ALL} réussissent : l'utilisateur existe dans tous les conteneurs.
  \item De retour dans \code{FREEPDB1}, \code{CREATE USER c##test ...} échoue avec \code{ORA-65094:
        invalid local user or role name} : un nom réservé aux comptes communs ne peut pas être créé comme
        utilisateur local dans une PDB.
  \item \code{DBA\_USERS} : \code{ALI -> COMMON = NO} (local) et \code{C##SUPERVISEUR -> COMMON = YES}
        (commun).
\end{enumerate}

\ecart{Le code \code{ORA-65096} attendu sur \code{CREATE USER bob} dans la racine n'est pas apparu : la
commande a réussi et \code{BOB} a été créé comme utilisateur \emph{commun} (\code{COMMON = YES} dans
\code{CDB\_USERS}). En 26ai Free de ce conteneur, le paramètre \code{COMMON\_USER\_PREFIX} est vide
(\code{SHOW PARAMETER common\_user\_prefix} ne renvoie aucune valeur), donc le préfixe \code{C\#\#}
n'est plus imposé pour créer un compte dans \code{CDB\$ROOT}. Le second contrôle (\code{ORA-65094}) est,
lui, bien obtenu.}

\subsection{Exercice 1.6 -- Supprimer un utilisateur}
\textbf{Objectif.} Supprimer un utilisateur qui possède des objets puis gérer une session ouverte.

\shots{1_6}
\shots{1_6_a}
\shots{1_6_drop}

\textbf{Observations et réponses.}
\begin{enumerate}
  \item \code{DROP USER ali} échoue avec \code{ORA-01922: CASCADE must be specified to drop 'ALI'} car ALI
        possède la table \code{NOTES} et un quota sur \code{USERS}.
  \item Une session ALI est maintenue ouverte en parallèle (script \code{1\_6\_a}, \code{SLEEP(120)}). Le
        \code{DROP USER ali CASCADE} échoue alors avec \code{ORA-01940: cannot drop a user who is currently
        connected}. La session est identifiée dans \code{V\$SESSION} (\code{SID=221, SERIAL\#=47280}).
  \item La session est terminée par \code{ALTER SYSTEM KILL SESSION '221,47280' IMMEDIATE} (le bloc
        signale \code{ORA-00031: session marked for kill}); la session côté utilisateur reçoit
        \code{ORA-00028: session has been terminated}. Après disparition de la session, \code{DROP USER ali
        CASCADE} réussit.
  \item Après suppression : \code{tables\_ali = 0} et \code{quotas\_ali = 0} dans \code{DBA\_TABLES} et
        \code{DBA\_TS\_QUOTAS}. La table \code{NOTES} et le quota disparaissent donc avec l'utilisateur.
\end{enumerate}
}{}
\IfFileExists{report/ch2.tex}{\section{Privilèges}

\subsection{Exercice 2.1 -- Le développeur minimal}
\textbf{Objectif.} Donner à DEV1 exactement les privilèges nécessaires pour travailler dans son propre
schéma, avec un quota borné.

\shots{2_1}

\textbf{Observations et réponses.}
\begin{enumerate}
  \item Les quatre privilèges accordés (\code{CREATE TABLE}, \code{CREATE VIEW}, \code{CREATE SEQUENCE},
        \code{CREATE PROCEDURE}) plus le quota \code{20M} sur \code{USERS} suffisent. Aucun rôle et aucun
        privilège \code{ANY} n'ont été donnés.
  \item \code{SESSION\_PRIVS} en DEV1 liste exactement : \code{CREATE PROCEDURE}, \code{CREATE SEQUENCE},
        \code{CREATE SESSION}, \code{CREATE TABLE} et \code{CREATE VIEW}. La table \code{PROJETS}, la
        séquence \code{SEQ\_PROJETS} et la vue \code{V\_PROJETS} sont créées avec succès.
  \item Il ne faut pas donner \code{UNLIMITED TABLESPACE} : ce privilège système accorde un espace
        illimité dans \emph{tous} les tablespaces et court-circuite le contrôle de quota. Un quota de 20 Mo
        sur \code{USERS} borne au contraire la consommation et évite qu'un développeur remplisse le disque.
\end{enumerate}

\subsection{Exercice 2.2 -- Lire sans modifier}
\textbf{Objectif.} Comparer les codes d'erreur d'un manque de droit DML et d'une table invisible, puis
distinguer \code{SELECT} de \code{READ}.

\shots{2_2}

\textbf{Observations et réponses.}
\begin{enumerate}
  \item En LECTEUR1, \code{SELECT COUNT(*)} sur \code{rh.employes} renvoie 8 : la lecture est autorisée.
  \item \code{UPDATE rh.employes SET salaire = 1} échoue avec \code{ORA-41900: missing UPDATE privilege on
        "RH"."EMPLOYES"} : la table est visible mais le droit d'écriture manque.
  \item \code{SELECT * FROM rh.departements} échoue avec \code{ORA-00942: table or view "RH"."DEPARTEMENTS"
        does not exist} : sans aucun privilège, la table est \emph{invisible} pour l'utilisateur, comme si
        elle n'existait pas. La différence est donc : \code{ORA-41900} = objet connu mais droit insuffisant,
        \code{ORA-00942} = objet non visible faute de privilège.
  \item Après remplacement de \code{SELECT} par \code{READ}, le \code{SELECT} fonctionne toujours, mais
        \code{SELECT ... FOR UPDATE} échoue avec \code{ORA-41900: missing SELECT privilege} : \code{READ}
        n'autorise que la lecture pure et \textbf{pas le verrouillage} des lignes.
\end{enumerate}

\ecart{Dans cette version, un manque de droit DML renvoie \code{ORA-41900} (« missing UPDATE/SELECT
privilege ») là où l'énoncé 19c attend \code{ORA-01031} (« insufficient privileges »). \code{ORA-00942}
reste identique.}

\subsection{Exercice 2.3 -- Droits par colonne, SQL92\_SECURITY et vue}
\textbf{Objectif.} Autoriser la modification de deux colonnes seulement, comprendre l'effet de
\code{SQL92\_SECURITY}, puis passer par une vue pour ne jamais exposer les salaires.

\shots{2_3}
\shots{2_3b}

\textbf{Observations et réponses.}
\begin{enumerate}
  \item Avec \code{GRANT UPDATE (poste, dept\_id) ON rh.employes TO omar}, la modification du salaire
        échoue : \code{ORA-41900: missing UPDATE privilege}.
  \item \code{UPDATE rh.employes SET poste = 'Architecte' WHERE emp\_id = 3} échoue aussi avec
        \code{ORA-41900: missing SELECT privilege} alors qu'elle ne touche qu'une colonne autorisée :
        \code{SHOW PARAMETER sql92\_security} renvoie \code{TRUE}. Avec \code{SQL92\_SECURITY = TRUE}, un
        \code{UPDATE}/\code{DELETE} comportant une clause \code{WHERE} exige en plus le droit
        \code{SELECT} sur les colonnes lues dans le \code{WHERE} (ici \code{emp\_id}).
  \item Sans clause \code{WHERE}, \code{UPDATE rh.employes SET poste = 'Architecte'} réussit
        (\code{8 rows updated}) : aucune colonne n'est lue pour filtrer. Puis \code{ROLLBACK}.
  \item Solution : accorder \code{SELECT, UPDATE (poste, dept\_id)} sur la vue \code{rh.v\_annuaire}
        (qui ne contient pas le salaire), puis retirer les droits sur la table. La modification via la vue
        réussit (\code{1 row updated}) et la tentative sur \code{salaire} via la vue échoue avec
        \code{ORA-00904: "SALAIRE": invalid identifier}, la colonne n'existant pas dans la vue.
  \item État final (\code{DBA\_COL\_PRIVS} / \code{DBA\_TAB\_PRIVS}) : OMAR ne possède plus que
        \code{UPDATE} sur \code{V\_ANNUAIRE.POSTE} et \code{V\_ANNUAIRE.DEPT\_ID} (privilèges de colonne) et
        \code{SELECT} sur \code{RH.V\_ANNUAIRE}. Toute modification directe de \code{RH.EMPLOYES} renvoie
        \code{ORA-00942} (table redevenue invisible).
\end{enumerate}

\ecart{Le \code{REVOKE UPDATE (poste, dept\_id) ON rh.employes FROM omar} est refusé en 26ai :
\code{ORA-01750: UPDATE/REFERENCES may only be REVOKEd from the whole table, not by column}. La
révocation a donc été faite au niveau table (\code{REVOKE UPDATE ON rh.employes FROM omar}), ce qui
supprime bien les deux privilèges de colonne.}

\subsection{Exercice 2.4 -- GRANT OPTION : la cascade}
\textbf{Objectif.} Observer la propagation d'un privilège objet avec \code{WITH GRANT OPTION} et sa
révocation en cascade.

\shots{2_4}

\textbf{Observations et réponses.}
\begin{enumerate}
  \item \code{DBA\_TAB\_PRIVS} sur \code{RH.DEPARTEMENTS} montre deux lignes :
        \code{A1 / GRANTOR=RH / GRANTABLE=YES} et \code{A2 / GRANTOR=A1 / GRANTABLE=NO}. On voit donc qui a
        accordé quoi et si le bénéficiaire peut le retransmettre.
  \item Après \code{REVOKE SELECT ON rh.departements FROM a1}, A2 perd l'accès : sa requête renvoie
        \code{ORA-00942}. La révocation de \code{GRANT OPTION} se propage \textbf{en cascade} aux
        privilèges dérivés.
\end{enumerate}

\subsection{Exercice 2.5 -- ADMIN OPTION : pas de cascade}
\textbf{Objectif.} Montrer que la révocation d'un privilège système avec \code{ADMIN OPTION} ne retire pas
les privilèges que le bénéficiaire a lui-même accordés.

\shots{2_5}

\textbf{Observations et réponses.}
\begin{enumerate}
  \item A1 accorde \code{CREATE VIEW} à A2 grâce à son \code{ADMIN OPTION}.
  \item Après \code{REVOKE CREATE VIEW FROM a1}, A2 possède \emph{toujours} \code{CREATE VIEW} :
        \code{USER\_SYS\_PRIVS} le renvoie et \code{DBA\_SYS\_PRIVS} contient seulement la ligne
        \code{A2 / CREATE VIEW / NO}. Il n'y a \textbf{pas de cascade} pour les privilèges système.
\end{enumerate}

\begin{center}
\small
\begin{tabular}{|p{3.2cm}|p{4.3cm}|p{4.3cm}|}
\hline
 & \textbf{GRANT OPTION} (objet) & \textbf{ADMIN OPTION} (système) \\
\hline
Transmission & possible à un autre utilisateur & possible à un autre utilisateur \\
\hline
Révocation du grantor & \textbf{cascade} : les droits dérivés tombent & \textbf{pas de cascade} : les droits dérivés restent \\
\hline
Traçabilité & \code{GRANTOR} visible dans \code{DBA\_TAB\_PRIVS} & droit « orphelin » plus difficile à retrouver \\
\hline
\end{tabular}
\end{center}

\textbf{La plus dangereuse pour un auditeur :} l'\code{ADMIN OPTION}, car la révocation ne se propage pas.
Les privilèges accordés en aval subsistent après la révocation du grantor et deviennent des droits
« orphelins » difficiles à repérer, alors qu'une révocation de \code{GRANT OPTION} nettoie automatiquement
toute la chaîne.

\subsection{Exercice 2.6 -- Le danger des privilèges ANY}
\textbf{Objectif.} Mesurer la portée d'un privilège \code{ANY} et lister les comptes à risque.

\shots{2_6}

\textbf{Observations et réponses.}
\begin{enumerate}
  \item Avec \code{SELECT ANY TABLE}, A2 lit \code{rh.employes} (8 lignes) et \code{dev1.projets}
        (0 ligne) sans aucun grant explicite : le privilège couvre toutes les tables de tous les schémas.
  \item En revanche \code{SELECT * FROM sys.user\$} échoue avec \code{ORA-00942} : le dictionnaire de
        données reste protégé (\code{O7\_DICTIONARY\_ACCESSIBILITY = FALSE}), \code{ANY} ne donne pas accès
        aux tables système.
  \item La jointure \code{DBA\_SYS\_PRIVS}/\code{DBA\_USERS} avec \code{ORACLE\_MAINTAINED = 'N'} et
        \code{PRIVILEGE LIKE '\%ANY\%'} ne renvoie qu'un compte : \code{A2} avec \code{SELECT ANY TABLE}.
        Le privilège est ensuite révoqué (\code{REVOKE SELECT ANY TABLE FROM a2}).
\end{enumerate}
}{}
\IfFileExists{report/ch3.tex}{\section{Rôles et modèle RBAC}

\subsection{Exercice 3.1 -- Construire la matrice RBAC du service RH}
\textbf{Objectif.} Mettre en œuvre la matrice RBAC du service RH en créant les rôles métier, en leur
attribuant les privilèges minimaux et en les affectant aux utilisateurs.

\begin{center}
\small
\begin{tabular}{|l|c|c|c|c|c|}
\hline
\textbf{Rôle / privilège} & \code{CREATE SESSION} & \code{rh.departements} & \code{rh.v_annuaire} & \code{rh.employes} (cols) & \code{rh.employes} (salaire) \\
\hline
\code{R_CONNEXION} & \textbf{Oui} & -- & -- & -- & -- \\
\hline
\code{R_RH_LECTEUR} & via \code{R_CONNEXION} & \textbf{SELECT} & \textbf{SELECT} & -- & -- \\
\hline
\code{R_RH_GESTION} & via \code{R_CONNEXION} & via \code{LECTEUR} & via \code{LECTEUR} & \textbf{SELECT, INSERT, UPD}(\code{poste},\code{dept_id}) & -- \\
\hline
\code{R_RH_PAIE} & via \code{R_CONNEXION} & -- & -- & \textbf{SELECT} & \textbf{UPD}(\code{salaire}) \\
\hline
\end{tabular}
\end{center}

\shots{3_1}

\textbf{Observations et réponses.}
\begin{enumerate}
  \item \textbf{LINA} (\code{R\_RH_LECTEUR}) : accède à \code{rh.departements} (4) et \code{rh.v_annuaire}
        (8), mais \code{SELECT * FROM rh.employes} renvoie \code{ORA-00942}.
  \item \textbf{OMAR} (\code{R\_RH_GESTION}) : accède à \code{rh.employes} (8), modifie \code{poste} avec
        succès (\code{1 row updated}), mais la modification de \code{salaire} renvoie \code{ORA-41900: missing
        UPDATE privilege}. Ses droits directs sur la table et la vue ont tous été révoqués au préalable
        (\code{DBA_TAB_PRIVS} et \code{DBA_COL_PRIVS} vides pour OMAR).
  \item \textbf{SARA} (\code{R\_RH_PAIE}) : lit \code{rh.employes}, modifie \code{salaire} avec succès
        (\code{1 row updated}), mais la modification de \code{poste} renvoie \code{ORA-41900} et l'accès à
        \code{rh.departements} renvoie \code{ORA-00942}.
\end{enumerate}

\subsection{Exercice 3.2 -- Lire la hiérarchie dans le dictionnaire}
\textbf{Objectif.} Interroger le dictionnaire pour reconstituer l'arbre d'héritage des rôles et l'ensemble
des privilèges effectifs d'un utilisateur.

\shots{3_2}

\textbf{Observations et réponses.}
\begin{enumerate}
  \item \code{DBA_ROLE_PRIVS WHERE grantee = 'OMAR'} ne renvoie que \code{R_RH_GESTION}. \code{R_RH_LECTEUR}
        et \code{R_CONNEXION} n'y apparaissent pas directement car ils sont accordés au rôle \code{R_RH_GESTION}
        et non directement à l'utilisateur OMAR.
  \item La requête hiérarchique avec \code{START WITH grantee = 'OMAR' CONNECT BY PRIOR granted_role = grantee}
        déroule l'arbre complet : \code{R_RH_GESTION} $\to$ \code{R_CONNEXION} et \code{R_RH_LECTEUR} $\to$
        \code{R_CONNEXION}.
  \item La jointure avec \code{DBA_TAB_PRIVS} / \code{DBA_COL_PRIVS} et l'interrogation de
        \code{ROLE_TAB_PRIVS} fournissent la liste exhaustive des 6 privilèges objet dont dispose OMAR
        via la hiérarchie des rôles.
\end{enumerate}

\subsection{Exercice 3.3 -- Effet immédiat ou à la reconnexion ?}
\textbf{Objectif.} Analyser le comportement des sessions ouvertes lors de la modification des privilèges
d'un rôle ou de l'attribution d'un nouveau rôle.

\shots{3_3_a}
\shots{3_3_b}
\shots{3_3_c}

\textbf{Observations et réponses.}
\begin{enumerate}
  \item \textbf{Étape 1} : quand l'administrateur accorde \code{SELECT ON rh.employes} au rôle
        \code{R_RH_LECTEUR}, la session de LINA déjà ouverte voit immédiatement la table (\code{8}) sans
        se reconnecter. Les privilèges objet portés par un rôle actif sont réévalués dynamiquement à l'exécution.
  \item \textbf{Étape 2} : après révocation du droit sur \code{R_RH_LECTEUR} et attribution du rôle
        \code{R_TEST} à LINA, le \code{SELECT} dans la même session échoue (\code{ORA-00942}). Le nouveau rôle
        accordé n'est pas actif dans une session déjà ouverte. Après \code{SET ROLE ALL}, le rôle s'active
        et la requête réussit (\code{8}). Après reconnexion, la requête réussit aussi.
  \item \textbf{Conclusion} : lors du départ d'un collaborateur, révoquer ses rôles ne suffit pas ; ses
        sessions ouvertes conservent leurs rôles actifs jusqu'à leur fermeture. Il faut obligatoirement
        \textbf{tuer les sessions ouvertes} avec \code{ALTER SYSTEM KILL SESSION} pour rendre la révocation
        immédiatement effective.
\end{enumerate}

\subsection{Exercice 3.4 -- Rôle protégé et rôles par défaut}
\textbf{Objectif.} Protéger le rôle de la paie par mot de passe et gérer les rôles par défaut.

\shots{3_4}

\textbf{Observations et réponses.}
\begin{enumerate}
  \item Après \code{ALTER ROLE r_rh_paie IDENTIFIED BY "Paie#2026"}, la tentative de connexion de Sara
        renvoie \code{ORA-01045: Login denied}. Le rôle protégé par mot de passe n'est plus activé
        automatiquement à la connexion ; comme Sara ne possédait \code{CREATE SESSION} que via
        \code{R_RH_PAIE}, elle ne peut plus ouvrir de session.
  \item Pour corriger, l'administrateur accorde \code{R_CONNEXION} directement à Sara et définit
        \code{ALTER USER sara DEFAULT ROLE r_connexion}. Sara se connecte alors avec le seul rôle
        \code{R_CONNEXION} actif ; l'accès à \code{rh.employes} renvoie \code{ORA-00942}.
  \item Sara tente \code{SET ROLE r_rh_paie IDENTIFIED BY "mauvais"} $\to$ \code{ORA-01979: missing or
        invalid password for role "R_RH_PAIE"}. Avec le bon mot de passe (\code{Paie#2026}), le rôle
        s'active (\code{Role set.}) et Sara peut à nouveau lire les employés et modifier leurs salaires.
\end{enumerate}

\ecart{Dans l'énoncé 19c, le rôle par défaut restait activé sans mot de passe à la connexion. En 26ai, un
rôle protégé par mot de passe n'est pas activé automatiquement s'il est le seul vecteur de connexion, ce qui
a provoqué \code{ORA-01045} jusqu'à l'octroi direct de \code{R_CONNEXION}.}

\subsection{Exercice 3.5 -- Le piège des rôles dans le code}
\textbf{Objectif.} Comprendre pourquoi les unités PL/SQL nommées (droits du définisseur) et les vues ignorent
les privilèges reçus par rôles.

\shots{3_5}

\textbf{Observations et réponses.}
\begin{enumerate}
  \item Bien qu'OMAR puisse exécuter \code{SELECT * FROM rh.employes} en SQL direct, la compilation de la
        fonction \code{nb_employes} échoue avec \code{PLS- / ORA-00942: table or view "RH"."EMPLOYES" does
        not exist}. De même, la création de la vue \code{v_effectif} renvoie \code{ORA-41904: Missing object
        privilege on "RH"."EMPLOYES" to create a view}.
  \item \textbf{Raison} : en PL/SQL à droits du définisseur (défaut) et pour la création de vues, le SGBD
        désactive les rôles lors de la compilation pour garantir la stabilité et la sécurité du code.
        Seuls les privilèges accordés \textbf{directement} à l'utilisateur sont pris en compte.
  \item L'option \code{AUTHID CURRENT_USER} a aussi échoué à la compilation en 26ai : le contrôle du SQL
        statique à la compilation exige que l'objet soit accessible au définisseur.
  \item \textbf{Correction la plus restreinte} : \code{GRANT SELECT ON rh.employes TO omar}. Ce privilège
        direct permet à la fonction de compiler sans erreur (\code{Function created.}) et à la vue d'être
        créée (\code{View created.}). C'est une exception assumée au principe RBAC pur.
\end{enumerate}

\subsection{Exercice 3.6 -- Que contiennent CONNECT et RESOURCE ?}
\textbf{Objectif.} Examiner le contenu des rôles système prédéfinis \code{CONNECT} et \code{RESOURCE}.

\shots{3_6}

\textbf{Observations et réponses.}
\begin{enumerate}
  \item \code{ROLE_SYS_PRIVS} montre : \code{CONNECT} contient \code{CREATE SESSION} et \code{SET CONTAINER}
        (2 privilèges). \code{RESOURCE} contient 15 privilèges de création (\code{CREATE TABLE}, \code{CREATE
        PROCEDURE}, \code{CREATE SEQUENCE}, \code{CREATE TRIGGER}, \code{CREATE TYPE}, etc.), mais
        \textbf{pas de quota} sur aucun tablespace.
  \item L'utilisateur \code{TEST_RES} doté de \code{CONNECT, RESOURCE} crée sa table \code{t_res}, mais
        l'insertion d'une ligne échoue avec \code{ORA-01950: insufficient quota on tablespace USERS}. Il
        lui manque un quota explicite (\code{UNLIMITED TABLESPACE} ayant été retiré de \code{RESOURCE}).
  \item Ces rôles historiques ne doivent pas servir de base à une politique RBAC d'entreprise : ils sont
        trop larges, figés, accordent des droits inutiles (ex: \code{CREATE CLUSTER}, \code{CREATE OPERATOR})
        et varient selon la version de la base sans correspondre aux métiers de l'organisation.
\end{enumerate}

\subsection{Exercice 3.7 -- Analyse des privilèges (DBMS\_PRIVILEGE\_CAPTURE)}
\textbf{Objectif.} Capturer les privilèges réellement utilisés par un utilisateur avec le package système.

\shots{3_7}

\textbf{Observations et réponses.}
\begin{enumerate}
  \item La capture \code{cap_omar} (\code{type => G_CONTEXT}, \code{condition => 'SESSION_USER = OMAR'})
        s'est créée (\code{CREATE_CAPTURE}), activée (\code{ENABLE_CAPTURE}) et désactivée
        (\code{DISABLE_CAPTURE}) sans erreur.
  \item Après \code{GENERATE_RESULT}, la consultation des vues \code{DBA_USED_OBJPRIVS} et
        \code{DBA_UNUSED_OBJPRIVS} est fonctionnelle (colonnes \code{CAPTURE}, \code{USERNAME},
        \code{USED_ROLE}, \code{OBJ_PRIV}, \code{OBJECT_NAME}).
  \item Cet outil permet en production d'analyser l'activité réelle d'un compte sur une période donnée pour
        retirer les privilèges inutilisés et appliquer rigoureusement le principe du moindre privilège.
\end{enumerate}
}{}
\IfFileExists{report/ch4.tex}{\section{Profils}

\subsection{Exercice 4.1 -- Profil de ressources}
\textbf{Objectif.} Vérifier \code{RESOURCE\_LIMIT}, créer le profil \code{P\_RH\_TEST}
(\code{SESSIONS\_PER\_USER 2}, \code{IDLE\_TIME 1}) et l'attribuer à LINA, puis observer
les limites de sessions simultanées et d'inactivité.

\shots{4_1_setup}
\shots{4_1_a}
\shots{4_1_b}
\shots{4_1_third}
\shots{4_1_idle}
\shots{4_1_idle_status}

\textbf{Observations et réponses.}
\begin{enumerate}
  \item \code{RESOURCE\_LIMIT} vaut \code{TRUE} : les limites de profil sont appliquées.
  \item Le profil \code{P\_RH\_TEST} est créé avec \code{SESSIONS\_PER\_USER = 2} et
        \code{IDLE\_TIME = 1} (1 minute), puis attribué à LINA.
  \item La troisième session LINA renvoie \code{ORA-02391: exceeded simultaneous
        SESSIONS\_PER\_USER limit} : le plafond de 2 sessions simultanées est respecté.
  \item Après \code{HOST sleep 150} (> IDLE\_TIME 1 minute), la requête renvoie
        \code{ORA-02396: exceeded maximum idle time, please connect again}. Pendant
        l'attente, le DBA voit la session LINA en statut \code{ACTIVE} dans \code{V\$SESSION},
        puis elle disparaît (\emph{sniped} par Oracle).
  \item Dans \code{DBA\_PROFILES}, la valeur \code{DEFAULT} dans la colonne \code{LIMIT}
        signifie que le paramètre reprend la valeur définie dans le profil \code{DEFAULT}.
\end{enumerate}
\ecart{Aucun. Le comportement est conforme à l'énoncé.}

\subsection{Exercice 4.2 -- Verrouillage après échecs}
\textbf{Objectif.} Configurer le verrouillage de compte après échecs de connexion
(\code{FAILED\_LOGIN\_ATTEMPTS 3}, \code{PASSWORD\_LOCK\_TIME 2/1440}) et en observer
l'effet ainsi que le déverrouillage automatique.

\shots{4_2}

\textbf{Observations et réponses.}
\begin{enumerate}
  \item Les trois premières connexions LINA avec un mauvais mot de passe renvoient
        \code{ORA-01017}. La quatrième, avec le bon mot de passe, renvoie
        \code{ORA-28000: The account is locked; login denied.} Le compte est verrouillé.
  \item \code{DBA\_USERS} affiche \code{ACCOUNT\_STATUS = LOCKED(TIMED)} avec un
        \code{LOCK\_DATE} renseigné.
  \item Après une attente de 130 secondes (> 2 minutes de \code{PASSWORD\_LOCK\_TIME}),
        la reconnexion de LINA réussit : le déverrouillage est automatique.
\end{enumerate}
\ecart{Aucun. Conforme à l'énoncé.}

\subsection{Exercice 4.3 -- Expiration, grâce et réutilisation}
\textbf{Objectif.} Configurer la réutilisation des mots de passe
(\code{PASSWORD\_REUSE\_TIME 30}, \code{PASSWORD\_REUSE\_MAX 3}), puis tester l'expiration
et la période de grâce avec \code{PASSWORD\_LIFE\_TIME} et \code{PASSWORD\_GRACE\_TIME}
réglés à 1 minute.

\shots{4_3}

\textbf{Observations et réponses.}
\begin{enumerate}
  \item Les règles \code{PASSWORD\_REUSE\_TIME = 30} et \code{PASSWORD\_REUSE\_MAX = 3}
        sont appliquées à \code{P\_RH\_TEST}.
  \item LINA change son mot de passe, puis tente de reprendre l'ancien immédiatement :
        la reprise a \textbf{réussi} (pas de \code{ORA-28007}) : l'historique de
        réutilisation n'a pas bloqué la réutilisation immédiate dans ce contexte.
        \ecart{L'énoncé 19c attend un refus (ORA-28007) en réutilisant l'ancien mot de
        passe ; en 26ai la reprise a fonctionné. Cela est consigné tel quel dans le log.}
  \item Avec \code{PASSWORD\_LIFE\_TIME = 1/1440} et \code{PASSWORD\_GRACE\_TIME =
        1/1440} : après 1 minute, la connexion est acceptée avec l'avertissement
        \code{ORA-28098} et \code{ACCOUNT\_STATUS = EXPIRED(GRACE)}. Après 2 minutes
        supplémentaires, la connexion renvoie \code{ORA-28001} et le compte passe en
        \code{EXPIRED}.
  \item Restauration à \code{PASSWORD\_LIFE\_TIME 90} / \code{PASSWORD\_GRACE\_TIME 7} :
        LINA reçoit un mot de passe neuf (\code{Lina\#2032}) et repasse en \code{OPEN}.
\end{enumerate}

\subsection{Exercice 4.4 -- Complexité des mots de passe}
\textbf{Objectif.} Associer une fonction de vérification de complexité
(\code{ORA12C\_VERIFY\_FUNCTION}) au profil et observer le rejet des mots de passe faibles
ainsi que l'obligation de la clause \code{REPLACE}.

\shots{4_4}

\textbf{Observations et réponses.}
\begin{enumerate}
  \item \code{ORA12C\_VERIFY\_FUNCTION} existe bien (objet \code{FUNCTION} propriétaire
        \code{SYS}). Elle est associée à \code{P\_RH\_TEST}.
  \item Tentatives de mots de passe faibles (\code{lina}, \code{abcdefgh},
        \code{Lina2026x}) sans clause \code{REPLACE} : toutes renvoient
        \code{ORA-28221: REPLACE clause not specified}. Quand une fonction de vérification
        est définie, Oracle exige la clause \code{REPLACE} avant d'évaluer la complexité.
  \item Le mot de passe \code{Lina\#2033} (contient le nom d'utilisateur) renvoie
        \code{ORA-28003: password contains the user name}, même avec la clause
        \code{REPLACE}.
\end{enumerate}
\ecart{L'énoncé 19c attend \code{ORA-28003} pour les mots de passe faibles. En 26ai, sans
clause \code{REPLACE}, c'est \code{ORA-28221} qui est renvoyé en premier (exigence de la
clause REPLACE liée à la fonction de vérification). Le mot de passe de LINA reste
\code{Lina\#2032} (compte \code{OPEN}).}

\subsection{Exercice 4.5 -- Profil de production et synthèse}
\textbf{Objectif.} Créer un profil de production \code{P\_RH\_STANDARD}, le déployer sur
LINA, OMAR et SARA, supprimer \code{P\_RH\_TEST} avec \code{CASCADE} et observer le sort
des utilisateurs concernés.

\shots{4_5}

\textbf{Observations et réponses.}
\begin{enumerate}
  \item Le profil \code{P\_RH\_STANDARD} est créé avec des limites de production
        (\code{SESSIONS\_PER\_USER 2}, \code{IDLE\_TIME 30}, \code{FAILED\_LOGIN\_ATTEMPTS
        3}, \code{PASSWORD\_LIFE\_TIME 90}, \code{PASSWORD\_REUSE\_TIME 365},
        \code{PASSWORD\_VERIFY\_FUNCTION ORA12C\_VERIFY\_FUNCTION}, etc.).
  \item \code{DROP PROFILE p\_rh\_test CASCADE} supprime le profil et réaffecte
        automatiquement LINA, OMAR et SARA au profil \code{DEFAULT} (visible dans
        \code{DBA\_USERS}). Un utilisateur dont le profil est supprimé avec
        \code{CASCADE} retombe donc sur le profil \code{DEFAULT}.
  \item Les trois utilisateurs sont ensuite passés sur \code{P\_RH\_STANDARD} ; leur
        \code{ACCOUNT\_STATUS} reste \code{OPEN}.
\end{enumerate}
\ecart{Aucun. Conforme à l'énoncé.}
}{}
\IfFileExists{report/ch5.tex}{\section{Audit unifié et FGA}

\subsection{Exercice 5.1 -- État de l'audit}
\textbf{Objectif.} Vérifier le mode d'audit unifié, lister les politiques prédéfinies et celles qui sont activées.

\shots{5_1}

\textbf{Observations et réponses.}
\begin{enumerate}
  \item \code{V\$OPTION} renvoie \code{Unified Auditing = TRUE} : la base est en mode audit unifié « pur » (26ai).
  \item \code{AUDIT\_UNIFIED\_ENABLED\_POLICIES} ne renvoie aucune ligne : aucune politique n'est activée par défaut.
  \item \code{COUNT(DISTINCT policy\_name) FROM audit\_unified\_policies} renvoie 16 politiques prédéfinies, parmi lesquelles
        \code{ORA\_LOGON\_FAILURES}, \code{ORA\_ACCOUNT\_MGMT}, \code{ORA\_SECURECONFIG}, \code{ORA\_STIG\_RECOMMENDATIONS},
        \code{ORA\_DV\_SCHEMA\_CHANGES}, etc.
\end{enumerate}
\ecart{Aucun. L'état initial (aucune politique activée) est conforme à l'observation sur 26ai.}

\subsection{Exercice 5.2 -- Échecs de connexion}
\textbf{Objectif.} Activer la politique \code{ORA\_LOGON\_FAILURES}, provoquer des échecs de connexion (OMAR ×2, PIRATE) et
en retrouver la trace dans \code{UNIFIED\_AUDIT\_TRAIL}.

\shots{5_2}

\textbf{Observations et réponses.}
\begin{enumerate}
  \item \code{AUDIT POLICY ora\_logon\_failures} : \code{Audit succeeded.}
  \item Les connexions OMAR/\code{faux} (×2) et PIRATE/\code{x} renvoient \code{ORA-01017: invalid credential or not authorized; logon denied.}
  \item Dans \code{UNIFIED\_AUDIT\_TRAIL} (\code{action\_name = 'LOGON' AND return\_code <> 0}) : trois enregistrements
        (\code{OMAR}, \code{OMAR}, \code{PIRATE}) avec \code{unified\_audit\_policies = ORA\_LOGON\_FAILURES}.
        Le code retour d'un mauvais mot de passe est \code{1017}.
\end{enumerate}
\ecart{Deux écarts 26ai : (1) \code{DBMS\_AUDIT\_MGMT.FLUSH\_UNIFIED\_AUDIT\_TRAIL} renvoie
\code{PLS-00302: component 'FLUSH\_UNIFIED\_AUDIT\_TRAIL' must be declared} (procédure non déclarée), mais les traces
sont néanmoins écrites directement dans \code{UNIFIED\_AUDIT\_TRAIL} ; (2)
\code{AUDIT\_UNIFIED\_ENABLED\_POLICIES} n'affiche pas la politique \code{ORA\_LOGON\_FAILURES} alors qu'elle est active
(les traces le prouvent).}

\subsection{Exercice 5.3 -- Tracer l'administration des comptes}
\textbf{Objectif.} Activer \code{ORA\_ACCOUNT\_MGMT}, créer l'utilisateur \code{STAGIAIRE}, lui accorder puis retirer
\code{R\_RH\_LECTEUR}, le supprimer, et reconstituer l'historique depuis la trace.

\shots{5_3}

\textbf{Observations et réponses.}
\begin{enumerate}
  \item \code{AUDIT POLICY ora\_account\_mgmt} : \code{Audit succeeded.}
  \item Les opérations sont toutes réussies : \code{CREATE USER stagiaire}, \code{GRANT CREATE SESSION},
        \code{GRANT r\_rh\_lecteur TO stagiaire}, \code{REVOKE r\_rh\_lecteur FROM stagiaire}, \code{DROP USER stagiaire}.
  \item La trace \code{UNIFIED\_AUDIT\_TRAIL} (\code{unified\_audit\_policies = ORA\_ACCOUNT\_MGMT}) contient les 5
        enregistrements (\code{CREATE USER}, \code{GRANT} ×2, \code{REVOKE}, \code{DROP USER}) avec \code{DBUSERNAME = SYS}.
  \item Le mot de passe n'apparaît pas en clair dans \code{SQL\_TEXT} : il est masqué par
        \code{CREATE USER stagiaire IDENTIFIED BY *}.
\end{enumerate}
\ecart{\code{DBMS\_AUDIT\_MGMT.FLUSH\_UNIFIED\_AUDIT\_TRAIL} renvoie \code{PLS-00302} (écart récurrent sur 26ai), mais les
traces sont présentes sans flush.}

\subsection{Exercice 5.4 -- Politique sur les données sensibles}
\textbf{Objectif.} Créer et activer une politique \code{pol\_rh\_salaires} sur \code{rh.employes} (\code{SELECT, UPDATE, DELETE}),
puis vérifier que les accès de SARA (paie), LINA (vue) et OMAR sont tracés.

\shots{5_4}

\textbf{Observations et réponses.}
\begin{enumerate}
  \item \code{CREATE AUDIT POLICY pol\_rh\_salaires ACTIONS SELECT, UPDATE, DELETE ON rh.employes} puis
        \code{AUDIT POLICY pol\_rh\_salaires} : \code{Audit succeeded.}
  \item SARA (\code{SET ROLE r\_rh\_paie IDENTIFIED BY "Paie\#2026"}) : \code{SELECT} sur \code{rh.employes} et
        \code{UPDATE rh.employes SET salaire = salaire * 1.02 WHERE emp\_id = 7} (\code{1 row updated}) sont tracés.
  \item LINA lit \code{rh.v\_annuaire} : deux enregistrements sont produits, l'un sur \code{OBJECT\_NAME = EMPLOYES}
        (table de base sous-jacente) et l'autre sur \code{OBJECT\_NAME = V\_ANNUAIRE} (la vue). L'accès via une vue est donc
        tracé sur la table de base \textbf{et} sur la vue.
  \item OMAR lit \code{rh.employes} : tracé (\code{OBJECT\_NAME = EMPLOYES}).
  \item \code{NOAUDIT POLICY pol\_rh\_salaires} puis \code{AUDIT POLICY pol\_rh\_salaires WHENEVER NOT SUCCESSFUL} :
        la politique passe en mode « échecs seulement ».
\end{enumerate}
\ecart{\code{AUDIT\_UNIFIED\_ENABLED\_POLICIES} n'affiche pas \code{pol\_rh\_salaires} après activation (écart d'affichage
récurrent). \code{FLUSH\_UNIFIED\_AUDIT\_TRAIL} renvoie \code{PLS-00302} (écart récurrent).}

\subsection{Exercice 5.5 -- Politique conditionnelle et rôle}
\textbf{Objectif.} Tester une politique conditionnelle avec clause \code{WHEN} sur \code{SYS\_CONTEXT}, puis auditer le rôle
\code{DBA} pour montrer qu'il couvre tous ses privilèges.

\shots{5_5}

\textbf{Observations et réponses.}
\begin{enumerate}
  \item \code{SYS\_CONTEXT('USERENV','CLIENT\_PROGRAM\_NAME')} renvoie
        \code{sqlplus@325ea919716d (TNS V1-V3)}.
  \item La création de \code{POL\_RH\_ADHOC} avec une clause \code{WHEN} contenant un littéral de chaîne renvoie
        \code{ORA-01756: quoted string not properly terminated} : la clause \code{WHEN} avec comparaison sur littéral texte
        ne passe pas en 26ai (la politique n'est pas créée).
  \item \code{CREATE AUDIT POLICY pol\_dba ROLES dba} puis \code{AUDIT POLICY pol\_dba} : \code{Audit succeeded.}
  \item \code{GRANT dba TO a1}, puis A1 active le rôle (\code{SET ROLE dba}) et crée \code{test\_dba} : l'action
        \code{CREATE USER} par A1 est tracée avec
        \code{unified\_audit\_policies = ORA\_ACCOUNT\_MGMT, POL\_DBA}. Auditer un rôle couvre donc ses privilèges
        (ici \code{CREATE USER}) d'un seul coup.
  \item \code{REVOKE dba FROM a1} et \code{DROP USER test\_dba} sont aussi tracés.
\end{enumerate}
\ecart{L'énoncé 19c attend une politique \code{POL\_RH\_ADHOC} avec clause \code{WHEN} sur
\code{SYS\_CONTEXT(...CLIENT\_PROGRAM\_NAME...)}. En 26ai, cette création renvoie \code{ORA-01756} (écart : la syntaxe de la
clause \code{WHEN} avec littéral texte n'est pas acceptée). \code{AUDIT\_UNIFIED\_ENABLED\_POLICIES} n'affiche pas les
politiques activées (écart récurrent).}

\subsection{Exercice 5.6 -- Audit fin des gros salaires (FGA)}
\textbf{Objectif.} Mettre en place une politique FGA sur \code{rh.employes} (\code{SALAIRE > 30000}, colonne \code{SALAIRE})
et vérifier qu'un seul accès de SARA est réellement audité.

\shots{5_6}

\textbf{Observations et réponses.}
\begin{enumerate}
  \item \code{DBMS\_FGA.ADD\_POLICY(... policy\_name => 'FGA\_GROS\_SALAIRES', audit\_condition => 'SALAIRE > 30000',
        audit\_column => 'SALAIRE', statement\_types => 'SELECT,UPDATE')} : proc\'edure termin\'ee.
  \item SARA (\code{SET ROLE r\_rh\_paie}) : (a) \code{SELECT emp\_id, nom} (pas de colonne \code{SALAIRE}) $\to$ pas de
        d\'eclenchement ; (b) \code{SELECT ... WHERE dept\_id = 40} (Ghazi 7650, Haddad 19000, aucun $> 30000$) $\to$ pas de
        d\'eclenchement ; (c) \code{SELECT ... WHERE dept\_id = 10} (Alaoui 45000 $> 30000$) $\to$ d\'eclenche l'audit FGA.
  \item \code{UNIFIED\_AUDIT\_TRAIL} (\code{fga\_policy\_name = FGA\_GROS\_SALAIRES}) contient \textbf{une seule ligne}
        (la requête de SARA sur le dept 10).
  \item \code{DBA\_FGA\_AUDIT\_TRAIL} renvoie \code{ORA-00904: "LSQL\_TEXT": invalid identifier} : en 26ai, la vue
        \code{DBA\_FGA\_AUDIT\_TRAIL} a une structure différente / ne contient pas les enregistrements (ils sont dans
        \code{UNIFIED\_AUDIT\_TRAIL}).
  \item \code{DBMS\_FGA.DISABLE\_POLICY('RH','EMPLOYES','FGA\_GROS\_SALAIRES')} : d\'esactiv\'ee.
\end{enumerate}
\ecart{Deux écarts 26ai : (1) \code{DBA\_FGA\_AUDIT\_TRAIL} n'est pas remplie / a des colonnes différentes (ORA-00904) :
les enregistrements FGA sont dans \code{UNIFIED\_AUDIT\_TRAIL} ; (2) \code{FLUSH\_UNIFIED\_AUDIT\_TRAIL} renvoie
\code{PLS-00302} (écart récurrent).}

\subsection{Exercice 5.7 -- Séparation des tâches et purge}
\textbf{Objectif.} Créer un utilisateur \code{AUDITEUR} (\code{CREATE SESSION} + \code{AUDIT\_VIEWER}), produire un rapport
d'audit, vérifier qu'il ne peut pas supprimer une politique, puis purger la trace unifiée.

\shots{5_7}

\textbf{Observations et réponses.}
\begin{enumerate}
  \item \code{CREATE USER auditeur}, \code{GRANT CREATE SESSION}, \code{GRANT AUDIT\_VIEWER TO auditeur} : cr\'eation r\'eussie.
  \item Rapport 24 h : \code{SELECT dbusername, action\_name, COUNT(*) ... GROUP BY dbusername, action\_name} sur
        \code{UNIFIED\_AUDIT\_TRAIL} (30 lignes : SYS, LINA, OMAR, SARA, A1, AUDITEUR, PIRATE).
  \item AUDITEUR tente \code{DROP AUDIT POLICY pol\_rh\_salaires} $\to$ \code{ORA-01031: insufficient privileges}. La
        s\'eparation des tâches est respectée : \code{AUDIT\_VIEWER} ne permet pas de supprimer une politique (il faut
        \code{AUDIT\_ADMIN}).
  \item Suppression des politiques de test (\code{NOAUDIT POLICY} puis \code{DROP AUDIT POLICY}). Le \code{DROP} sur une
        politique encore activée renvoie \code{ORA-46361} (première exécution) ; il faut d'abord \code{NOAUDIT}.
  \item Purge : \code{DBMS\_AUDIT\_MGMT.CLEAN\_AUDIT\_TRAIL(...)} termin\'ee. Après purge, la relecture de
        \code{UNIFIED\_AUDIT\_TRAIL} ne montre plus qu'une trace r\'esiduelle (la purge elle-même, \code{SYS EXECUTE
        ORA$MANDATORY}) : la trace a bien été vidée.
\end{enumerate}
\ecart{\code{DBMS\_AUDIT\_MGMT.FLUSH\_UNIFIED\_AUDIT\_TRAIL} renvoie \code{PLS-00302} (écart récurrent sur 26ai), mais la
purge \code{CLEAN\_AUDIT\_TRAIL} fonctionne.}
}{}
\IfFileExists{report/ch6.tex}{\section{Étude de cas : Banque Atlas}

\subsection{6\_init — Initialisation du schéma BANQUE}
\textbf{Objectif.} Créer le schéma BANQUE (tables clients, comptes, opérations, vue v\_clients\_guichet,
procédure p\_operation) et peupler les données de test.

\shots{6_init}

\textbf{Observations et réponses.}
\begin{enumerate}
  \item L'utilisateur \code{BANQUE} est créé avec \code{NO AUTHENTICATION} et
        \code{QUOTA UNLIMITED ON users} (schéma seulement, pas de connexion directe).
  \item Les tables \code{clients} (4 lignes), \code{comptes} (5 lignes) et
        \code{operations} sont créées et peuplées.
  \item La vue \code{v\_clients\_guichet} expose \code{client\_id, nom, cin, telephone}
        (pas de \code{revenu\_mensuel}, pas de \code{conseiller}).
  \item La procédure \code{p\_operation} se compile sans erreur : elle est la seule porte
        d'entrée pour les mouvements d'argent (droits du définisseur), contrôle le solde et
        journalise dans \code{operations} avec \code{SYS\_CONTEXT('USERENV','SESSION\_USER')}.
\end{enumerate}
\ecart{Aucun.}

\subsection{Exercice 6.1 — T1 : Matrice RBAC}
\textbf{Objectif.} Construire la matrice rôles × objets et justifier chaque case au regard des
exigences E1 à E3.

\begin{center}
\small
\begin{tabular}{|l|c|c|c|c|c|c|}
\hline
\textbf{Rôle} & \code{CREATE SESSION} & \code{clients} & \code{comptes} & \code{operations} & \code{v\_clients\_guichet} & \code{p\_operation} \\
\hline
\code{R\_BQ\_CONNEXION} & \textbf{Oui} & -- & -- & -- & -- & -- \\
\hline
\code{R\_BQ\_GUICHET} & via CONNEXION & -- & SELECT & -- & SELECT & EXECUTE \\
\hline
\code{R\_BQ\_CONSEILLER} & via GUICHET & SELECT, INSERT & INSERT & SELECT & SELECT (hérité) & EXECUTE (hérité) \\
\hline
\code{R\_BQ\_DIRECTEUR} & via CONSEILLER & (hérité) & (hérité) & (hérité) & (hérité) & EXECUTE (hérité) \\
\hline
\code{R\_BQ\_AUDIT} & \textbf{Oui} & SELECT & SELECT & SELECT & -- & -- \\
\hline
\code{R\_BQ\_BATCH} & via CONNEXION & -- & SELECT & -- & -- & EXECUTE \\
\hline
\end{tabular}
\end{center}

\textbf{Héritages :} \code{R\_BQ\_GUICHET} $\to$ \code{R\_BQ\_CONNEXION} ;
\code{R\_BQ\_CONSEILLER} $\to$ \code{R\_BQ\_GUICHET} $\to$ \code{R\_BQ\_CONNEXION} ;
\code{R\_BQ\_DIRECTEUR} $\to$ \code{R\_BQ\_CONSEILLER} ; \code{R\_BQ\_AUDIT} et
\code{R\_BQ\_BATCH} $\to$ \code{R\_BQ\_CONNEXION}.

\textbf{Justification (E1–E3) :}
\begin{enumerate}
  \item \textbf{E1} (aucun privilège direct aux utilisateurs métier) : tous les privilèges
        sont accordés aux rôles, jamais directement aux comptes GUI1, GUI2, CONS1, DIR1,
        AUD1, APP\_BATCH.
  \item \textbf{E2} (pas de ANY, pas de DBA/CONNECT/RESOURCE direct) : aucun compte métier
        ne possède ces rôles/privilèges directement.
  \item \textbf{E3} (soldes modifiables uniquement par P\_OPERATION) : aucun rôle n'a
        \code{UPDATE} sur \code{comptes.solde} ; seule la procédure \code{p\_operation}
        (droits du définisseur) modifie les soldes.
\end{enumerate}
\textbf{Seul privilège avec retransmission :} \code{GRANT r\_bq\_guichet TO dir1 WITH ADMIN
OPTION} (accordé directement au compte DIR1, pas via le rôle R\_BQ\_DIRECTEUR).

\subsection{Exercice 6.2 — T2 : Profils}
\textbf{Objectif.} Créer \code{P\_BQ\_AGENT}, \code{P\_BQ\_CONSEIL}, \code{P\_BQ\_SENSIBLE}
et \code{P\_BQ\_SERVICE} conformément à E4 et E5.

\shots{6_2}

\textbf{Observations et réponses.}
\begin{enumerate}
  \item \code{P\_BQ\_AGENT} : \code{SESSIONS\_PER\_USER 1}, \code{IDLE\_TIME 10},
        \code{FAILED\_LOGIN\_ATTEMPTS 3}, \code{PASSWORD\_LOCK\_TIME 30/1440},
        \code{PASSWORD\_LIFE\_TIME 60}, \code{PASSWORD\_REUSE\_TIME 365},
        \code{PASSWORD\_REUSE\_MAX 5}, \code{PASSWORD\_VERIFY\_FUNCTION ORA12C\_VERIFY\_FUNCTION}.
  \item \code{P\_BQ\_CONSEIL} : identique à \code{P\_BQ\_AGENT} sauf
        \code{SESSIONS\_PER\_USER 2} (le conseiller a besoin de 2 sessions).
  \item \code{P\_BQ\_SENSIBLE} : mêmes règles E4 que l'agent, avec
        \code{ORA12C\_STRONG\_VERIFY\_FUNCTION} (E5 : mot de passe renforcé pour directeur
        et auditeur).
  \item \code{P\_BQ\_SERVICE} : \code{PASSWORD\_LIFE\_TIME UNLIMITED} (pas d'expiration),
        \code{IDLE\_TIME DEFAULT}, \code{SESSIONS\_PER\_USER DEFAULT} (plusieurs sessions),
        fonction forte, verrouillage court (éviter déni de service).
\end{enumerate}
\textbf{Gestion de la différence guichet (1 session) / conseiller (2 sessions) :}
deux profils distincts (\code{P\_BQ\_AGENT} et \code{P\_BQ\_CONSEIL}) car
\code{SESSIONS\_PER\_USER} est un paramètre de profil. Alternative : un seul profil +
trigger de connexion, mais moins standard et plus complexe à maintenir.
\ecart{Aucun.}

\subsection{Exercice 6.3 — T3 : Comptes}
\textbf{Objectif.} Créer GUI1, GUI2, CONS1, DIR1, AUD1, APP\_BATCH, DEV\_BQ avec le bon
profil, aucun quota, mot de passe conforme à la fonction de vérification, et le proxy
DEV\_BQ $\to$ BANQUE.

\shots{6_3d}

\textbf{Observations et réponses.}
\begin{enumerate}
  \item \code{GUI1}, \code{GUI2} (\code{P\_BQ\_AGENT}) ; \code{CONS1} (\code{P\_BQ\_CONSEIL}) ;
        \code{DIR1}, \code{AUD1} (\code{P\_BQ\_SENSIBLE}) ; \code{APP\_BATCH}
        (\code{P\_BQ\_SERVICE}) ; \code{DEV\_BQ} (\code{DEFAULT}).
  \item Aucun quota n'est affecté (la requête \code{DBA\_TS\_QUOTAS} est vide pour ces
        comptes).
  \item Les mots de passe respectent les fonctions de vérification (ex. \code{ORA12C\_VERIFY\_FUNCTION}
        rejette les mots de passe contenant le nom d'utilisateur ; plusieurs tentatives ont
        été nécessaires pour trouver des mots de passe conformes).
  \item \code{ALTER USER banque GRANT CONNECT THROUGH dev\_bq} autorise le proxy.
  \item \code{DEV\_BQ} n'a pas besoin de quota : les objets qu'il crée (index, vues)
        appartiennent au schéma \code{BANQUE} (via la connexion proxy
        \code{dev\_bq[banque]}), et \code{BANQUE} a le quota \code{UNLIMITED}.
\end{enumerate}
\ecart{Aucun.}

\subsection{Exercice 6.4 — T4 : Rôles et habilitations}
\textbf{Objectif.} Implémenter la matrice T1, attribuer les rôles aux comptes, déléguer
\code{R\_BQ\_GUICHET} à DIR1 avec ADMIN OPTION, et vérifier E1/E2.

\shots{6_4}
\shots{6_4c}

\textbf{Observations et réponses.}
\begin{enumerate}
  \item Tous les rôles sont créés avec la bonne hiérarchie d'héritage.
  \item Les rôles sont affectés aux comptes : GUI1/GUI2 $\to$ \code{R\_BQ\_GUICHET},
        CONS1 $\to$ \code{R\_BQ\_CONSEILLER}, DIR1 $\to$ \code{R\_BQ\_DIRECTEUR},
        AUD1 $\to$ \code{R\_BQ\_AUDIT}, APP\_BATCH $\to$ \code{R\_BQ\_BATCH},
        DEV\_BQ $\to$ \code{R\_BQ\_CONNEXION}.
  \item \code{GRANT r\_bq\_guichet TO dir1 WITH ADMIN OPTION} : DIR1 peut déléguer
        \code{R\_BQ\_GUICHET} (test 10 réussit) mais pas \code{R\_BQ\_CONSEILLER}
        (test 11 : \code{ORA-01932}).
  \item Contrôle E1 : tous les privilèges objet passent par les rôles (requête 6\_4c),
        aucun privilège direct aux utilisateurs.
  \item Contrôle E2 : aucun privilège système direct (\code{DBA\_SYS\_PRIVS} vide pour les
        utilisateurs métier), aucun rôle DBA/CONNECT/RESOURCE directement affecté.
  \item Arbre des rôles de DIR1 :
        \code{R\_BQ\_DIRECTEUR} $\to$ \code{R\_BQ\_CONSEILLER} $\to$ \code{R\_BQ\_GUICHET}
        $\to$ \code{R\_BQ\_CONNEXION}.
\end{enumerate}
\ecart{Aucun.}

\subsection{Exercice 6.5 — T5 : Audit}
\textbf{Objectif.} Mettre en place l'audit répondant à E6 à E8 :
\code{POL\_BQ\_CLIENTS} (modifications clients), \code{POL\_BQ\_SOLDES} (échecs seulement
sur comptes/opérations), FGA \code{FGA\_BQ\_GROS\_MONTANTS} (insertions $\ge$ 50 000),
et \code{ORA\_ACCOUNT\_MGMT} (E8).

\shots{6_5}

\textbf{Observations et réponses.}
\begin{enumerate}
  \item \code{POL\_BQ\_CLIENTS} (\code{INSERT, UPDATE, DELETE ON banque.clients}) :
        succès et échecs tracés (E6).
  \item \code{POL\_BQ\_SOLDES} et \code{POL\_BQ\_SOLDES\_OPS} activés avec
        \code{WHENEVER NOT SUCCESSFUL} : seules les tentatives de modification directe
        des soldes échouées sont tracées (E6).
  \item \code{FGA\_BQ\_GROS\_MONTANTS} sur \code{banque.operations} (\code{INSERT},
        \code{MONTANT >= 50000}) : trace les gros montants avec l'agent (E7).
  \item \code{ORA\_ACCOUNT\_MGMT} activé : trace CREATE USER, GRANT, REVOKE, DROP USER
        (E8).
\end{enumerate}
\textbf{Pourquoi POL\_BQ\_SOLDES ne trace pas les succès :} les modifications légitimes
des soldes passent par \code{p\_operation} (droits du définisseur) et sont déjà
journalisées dans la table \code{operations}. Tracer les succès noierait la trace sous du
bruit ; seules les tentatives directes et non autorisées (échecs) méritent d'être
signalées.
\ecart{\code{AUDIT\_UNIFIED\_ENABLED\_POLICIES} n'affiche pas les politiques activées
(écart récurrent 26ai).}

\subsection{Exercice 6.6 — T6 : Campagne de tests (16 tests)}
\textbf{Objectif.} Exécuter les 16 tests, comparer attendu vs obtenu, et corriger les écarts.

\shots{6_6_a}
\shots{6_6_b}
\shots{6_6_fond}

\begin{center}
\small
\begin{tabular}{|c|l|l|l|l|}
\hline
\textbf{N°} & \textbf{Compte} & \textbf{Action} & \textbf{Attendu} & \textbf{Obtenu} \\
\hline
1 & GUI1 & lire v\_clients\_guichet & succès & succès (4 lignes) \\
\hline
2 & GUI1 & lire revenu\_mensuel dans clients & ORA-00942 & ORA-00942 \\
\hline
3 & GUI1 & dépôt 2000 sur 1003 via p\_operation & succès & succès (solde 6200) \\
\hline
4 & GUI1 & UPDATE comptes solde = 1e6 & ORA-01031 + trace & ORA-41900 \\
\hline
5 & GUI1 & retrait 60000 sur 1002 via p\_operation & succès + FGA & succès + trace FGA \\
\hline
6 & GUI1 & retrait 5000 sur 1005 & ORA-20002 & ORA-20002 \\
\hline
7 & GUI1 & ouvrir 2e session & ORA-02391 & ORA-02391 \\
\hline
8 & CONS1 & modifier téléphone client 2 & succès + trace & succès + trace POL\_BQ\_CLIENTS \\
\hline
9 & CONS1 & modifier CIN client 2 & ORA-01031 & ORA-41900 \\
\hline
10 & DIR1 & accorder R\_BQ\_GUICHET à GUI3 & succès + trace & succès + trace ORA\_ACCOUNT\_MGMT \\
\hline
11 & DIR1 & accorder R\_BQ\_CONSEILLER à GUI3 & ORA-01932 & ORA-01932 \\
\hline
12 & AUD1 & lire operations + trace & succès & succès (3 opérations) \\
\hline
13 & AUD1 & DELETE FROM banque.operations & ORA-01031 & ORA-41900 \\
\hline
14 & GUI2 & 3 mauvais mdp puis le bon & ORA-28000 & ORA-28000 \\
\hline
15 & DEV\_BQ & connexion dev\_bq[banque] + index & succès & succès \\
\hline
16 & APP\_BATCH & dépôt intérêts 1200 sur 1002 & succès & succès (solde 92400) \\
\hline
\end{tabular}
\end{center}

\textbf{Écarts et corrections :}
\begin{enumerate}
  \item Tests 4, 9, 13 : l'énoncé attend \code{ORA-01031} mais 26ai renvoie
        \code{ORA-41900: missing UPDATE/DELETE privilege} (écart connu, chapitre 2).
        Le comportement est équivalent (refus d'accès) ; aucun correctif nécessaire.
  \item Test 15 : initialement échouait avec \code{ORA-01045} car \code{BANQUE} n'avait
        pas \code{CREATE SESSION}. Corrigé par \code{GRANT CREATE SESSION TO banque}.
  \item Le test 7 (2e session GUI1) a nécessité une session de fond ouverte en parallèle
        (\code{6\_6\_fond}) pour saturer la limite \code{SESSIONS\_PER\_USER 1}.
\end{enumerate}

\subsection{Exercice 6.7 — T7 : Revue des habilitations}
\textbf{Objectif.} Produire trois requêtes de revue pour l'auditeur.

\shots{6_7}
\shots{6_7b}

\textbf{Observations et réponses.}
\begin{enumerate}
  \item Requête 1 (\code{LISTAGG} des rôles par utilisateur) : 8 utilisateurs métier,
        chacun avec son profil, son état (\code{OPEN}), sa date d'expiration et ses rôles
        directs concaténés.
  \item Requête 2 : tous les privilèges objet sur le schéma \code{BANQUE} sont accordés
        aux rôles (jamais directement aux utilisateurs), conformément à E1.
  \item Requête 3 : événements d'audit par politique et utilisateur sur 7 jours
        (\code{ORA\_ACCOUNT\_MGMT}, \code{ORA\_LOGON\_FAILURES}, \code{POL\_BQ\_CLIENTS},
        \code{POL\_BQ\_SOLDES\_OPS}, \code{ORA\$MANDATORY}).
\end{enumerate}
\ecart{La requête 2 utilise \code{TABLE\_NAME} (et non \code{OBJECT\_NAME}) car
\code{DBA\_TAB\_PRIVS} n'a pas de colonne \code{OBJECT\_NAME} en 26ai.}

\subsection{Exercice 6.8 — Questions de synthèse}
\textbf{Objectif.} Répondre aux 6 questions de synthèse en prose.

\subsubsection{1. Départ de GUI2}
GUI2 quitte la banque. Procédure complète :
\begin{enumerate}
  \item \textbf{Verrouiller immédiatement} le compte :
        \code{ALTER USER gui2 ACCOUNT LOCK;} (empêche toute nouvelle connexion).
  \item \textbf{Tuer ses sessions ouvertes} :
        \code{SELECT sid, serial\# FROM v\$session WHERE username='GUI2';} puis
        \code{ALTER SYSTEM KILL SESSION 'sid,serial\#' IMMEDIATE;}
        (comme au chapitre 3.3, la révocation des rôles ne suffit pas pour les sessions
        déjà ouvertes).
  \item \textbf{Révoquer ses rôles} : \code{REVOKE r\_bq\_guichet FROM gui2;}
  \item \textbf{Vérifier qu'il ne possède aucun objet} :
        \code{SELECT owner, object\_name FROM dba\_objects WHERE owner='GUI2';}
  \item \textbf{Consulter son historique d'audit} dans \code{UNIFIED\_AUDIT\_TRAIL}
        pour garder une trace de ses dernières actions.
  \item \textbf{Supprimer le compte plus tard} (après la durée de conservation légale)
        : \code{DROP USER gui2 CASCADE;}
\end{enumerate}
\textbf{Faut-il supprimer immédiatement ?} Non. On verrouille et on révoque
immédiatement pour couper l'accès, mais on conserve le compte (et ses objets/audit)
pendant la durée de conservation requise avant suppression définitive.

\subsubsection{2. Pourquoi P\_OPERATION protège mieux qu'un GRANT UPDATE(solde)}
\code{P\_OPERATION} (droits du définisseur) centralise : contrôle de validité du montant
(\code{> 0}), vérification du solde suffisant pour les retraits (\code{ORA-20002}),
mise à jour atomique du solde \textbf{et} journalisation dans \code{operations} (agent,
montant, type, date), le tout dans une transaction avec \code{COMMIT}. Un
\code{GRANT UPDATE(solde)} direct aux guichetiers permettrait de modifier le solde
arbitrairement sans contrôle ni journalisation, contournant E3.

\subsubsection{3. Risque de la délégation du directeur}
DIR1 a \code{R\_BQ\_GUICHET} avec \code{ADMIN OPTION} : il peut créer des habilitations
excessives ou frauduleuses (accorder le rôle à n'importe qui). L'audit
\code{ORA\_ACCOUNT\_MGMT} trace chaque \code{GRANT}/\code{REVOKE} (qui, quand, quel rôle).
L'\code{ADMIN OPTION} ne se révoque pas en cascade : si on retire le rôle à DIR1, les
comptes qu'il a habilités conservent leur rôle. Il faut donc auditer et réviser
régulièrement.

\subsubsection{4. Le DBA (SYS) peut lire les revenus — acceptable ?}
Non, ce n'est pas acceptable dans un contexte strict car le DBA a un accès total aux
données sensibles (\code{revenu\_mensuel}). Solutions Oracle au-delà du TP :
Oracle Database Vault (séparation des tâches DBA / sécurité), Transparent Data
Encryption (chiffrement des données), Data Redaction (masquage dynamique), Virtual
Private Database (filtrage ligne par ligne), Oracle Label Security (étiquettes de
confidentialité), Audit Vault and Database Firewall (audit centralisé).

\subsubsection{5. Limiter le conseiller à ses clients}
Créer une vue filtrée :
\code{CREATE VIEW banque.v\_clients\_conseiller AS SELECT * FROM banque.clients WHERE
conseiller = SYS\_CONTEXT('USERENV','SESSION\_USER');} (les valeurs de la colonne
\code{conseiller} sont des noms d'utilisateur). Ou mettre en place une politique VPD
(\code{DBMS\_RLS}) qui ajoute automatiquement le prédicate
\code{conseiller = SYS\_CONTEXT('USERENV','SESSION\_USER')} sur les requêtes de
\code{clients}.

\subsubsection{6. Deux améliorations de la politique de l'agence}
\begin{enumerate}
  \item \textbf{Authentification forte / centralisée} : remplacer l'authentification
        mot de passe par Oracle Database Vault ou une authentification multifactorielle,
        et centraliser la gestion des comptes (LDAP / Active Directory) pour éviter la
        prolifération de mots de passe locaux.
  \item \textbf{Chiffrement et révision périodique} : activer TDE sur les tables
        sensibles (clients, opérations), chiffrer les échanges réseau (TLS), et mettre en
        place une revue trimestrielle des habilitations avec l'analyse des privilèges
        (\code{DBMS\_PRIVILEGE\_CAPTURE}) pour retirer les droits inutilisés.
\end{enumerate}
\ecart{Aucun.}
}{}

\end{document}
